\noindent This systematized review assessed the comparability of conventional automatic and gravimetric methods for ambient PM10 and PM2.5 and the factors associated with measurement differences. Searches in Scopus, Web of Science, and PubMed were combined with a bounded citation search. Among 437 unique records, 84 had an available report for assessment and 66 were selected. Thirty-two reports were not retrieved, and six references could not be sufficiently identified. Selection and focused data extraction were assisted by artificial intelligence. Descriptive synthesis distinguished particle fractions, technologies, instrument configurations, and raw or corrected results, without meta-analysis. Differences varied with heating, semivolatile material, humidity, aerosol properties, concentration, and data processing. Heated oscillating microbalance instruments frequently measured lower concentrations than gravimetry, whereas modified systems also produced higher values where reference filters lost particulate material. Beta attenuation and optical methods showed responses dependent on measurement conditions. Local corrections improved results without demonstrating universal transferability. Gravimetry was also subject to measurement artifacts. Inferences are limited by the single-researcher review, missing reports, and heterogeneity. Comparability therefore requires assessing the combination of instrument, configuration, reference method, and environmental conditions.

\vspace{\baselineskip}
\noindent\textbf{Keywords:} particulate matter; gravimetry; continuous monitoring; comparability; relative humidity.
